Let us now look at some more advanced topics in the theory of finite fields. 


\bigskip


In Tables~\ref{tab:finite:field:activities:A}-\ref{tab:finite:field:activities:B}, 
a summary of 
finite field activities is shown.
\orbitertableofcommands{finite_field_activities_summary_1}
\orbitertableofcommands{finite_field_activities_summary_2}
		

\bigskip




A normal basis for a field extension $\bbF_{q^d}$ over $\bbF_q$ 
is a basis of $\bbF_{q^d}$ as vector space over $\bbF_q$ 
which consists of one cycle of the Frobenius 
automorphism of $\bbF_{q^d}$ over $\bbF_q.$
For instance, the command

\sourcecode{normal_basis_2_3}

computes a normal basis of $\bbF_8$ over $\bbF_2.$
Using the polynomial $X^{3} + X^{2} + 1,$
the normal basis in terms of the 
standard polynomial basis $1,X,X^2,\ldots$ 
is given by the columns of the matrix 
$$
\left[
\begin{array}{ccc}
1 &0 &0\\ 
1 &1 &0 \\
1 &0 &1 \\
\end{array}
\right].
$$
Reading the columns as cofficient 
vectors with respect to the standard basis, 
the normal basis is 
$$
b_1 = 1 + X + X^2, \quad b_2 = X, \quad b_3 = X^2.
$$
Let us apply the Frobenius mapping $\Phi$ 
to the elements of the normal bases:
\begin{align*}
b_1^\Phi &= (1+X+X^2)^2 = 1 + X^2 + X^4 = 1 + X^2 + X^3 + X = 1 + X + X^2 + X^2 + 1 = X = b_2,\\
b_2^\Phi &= X^2 = b_3, \\
b_3^\Phi &= X^4 = X^3+X = X^2+X+1 = b_1.
\end{align*}
Thus,
$$
b_1 \mapsto b_2 \mapsto b_3 \mapsto b_1
$$ 
under $\Phi,$ as required.

\bigskip

A field is a vector space over any of its subfields. 
Using a field basis, the elements of the large field can be 
identified with invertible matrices. 
So, for $\bbF_{q^r}$ over $\bbF_q$, and for $a \in \bbF_{q^r},$ 
we consider the $\bbF_q$-linear map 
$$
\bbF_{q^r} \rightarrow \bbF_{q^r} , x \mapsto ax.
$$
The following code computes the field reduction from $\bbF_{64}$ to $\bbF_8$. 
Elements in the small field are represented as colors.
The $(i,j)$-th block is the matrix of $a=i8+j$ in the chosen basis.

\sourcecode{F_64_over_F8_field_reduction}

The output is shown below: 

\orbiteroutput{GRAPHICS/WRAPPERS/field_reduction_F64_to_F8}

Note that the dimension of the vector space is 2, 
so the block matrices are $2 \times 2.$
Observe that $\bbF_{64}$ has many subfields. 
The reduction from $\bbF_{64}$ to $\bbF_4$ and to $\bbF_2$ is also possible:

\orbiteroutput{GRAPHICS/WRAPPERS/field_reduction_F64_to_F4_and_F2}

Here, the block matrices have size $3 \times 3$ 
when reducing to $\bbF_4$ (left picture) 
and $6 \times 6$ when reducing to $\bbF_2$.

\bigskip

The minimum polynomials associated with the $n$-th roots 
over $\bbF_q$ can be computed using the \verb'-nth_roots' command, which is a 
finite field activity. 
The activity is applied to the field $\bbF_q$ 
over which the $n$-th roots are defined.
The command constructs the field extension $\bbF_{q^m}$ 
where $m$ is the order of $q$ modulo $n$. 
This field extension contains the $n$-th roots of unity. 
Let $\alpha$ be a primitive element of $\bbF_{q^m}$ 
and let $\beta$ be a generator of the subgroup of $n$-th roots. 
Also, let $\gamma$ be the generator of the subgroup of $q-1$ th roots, 
which are the elements of the multiplicative group of $\bbF_q.$ 
The output lists the $n$-th roots first, generated by $\beta$. 
After that, the $q-1$th roots are shown, generated by $\gamma$. 
Finally, a table is produced which shows the irreducible polynomials 
over $\bbF_q$ associated with the $n$-th roots of unity.
For instance, the following command computes the minimum polynomials 
of all $21$st roots of unity over $\bbF_8$:

\sourcecode{F_8_Nth_roots_21}

The output is:

\orbiteroutput{GRAPHICS/WRAPPERS/nth_roots_21_over_F8}


\bigskip

In Section~\ref{sec:finite:fields}, 
we have considered the Vandermonde matrix over $\bbF_7$. 
Let us do the same for the field $\bbF_8$ instead. 
We use the following command:

\sourcecode{F_8_vandermonde}

The output is shown below. 
Again, the first matrix is $V = (x_{i}^j).$ The second matrix is $V^{-1}$:

\orbiteroutput{GRAPHICS/WRAPPERS/F8_vandermonde}

\bigskip


Let us now do a somewhat larger example of the same problem. 
The next command computes the Vandermonde matrix and its inverse over the field 
$\bbF_{512}$:

\sourcecode{F_512_vandermonde}

The matrix is not shown. It would be too big to be printed.
In order to save disc space, 
we delete the output files, using the \verb'rm' command.

\bigskip






Orbiter can create code for the number theoretic transform. 
This is the discrete Fourier transform performed over finite fields.
The generated code can be compiled with the Orbiter library. 
Compiling code requires additional makefile options are necessary. 
Because of this, we define the following makefile variables 
at the top of the makefile.

\sourcecodevariable{SRC}


\bigskip

Suppose we want to create 
the number theoretic transform for the 16th 
roots of unity inside the field $\bbF_{17}.$ 
Here is the command to generate the Orbiter source code:

\sourcecode{NTT_k4_q17.cpp}

This produces a C++ file \verb'NTT_k4_q17.cpp'. 
This file should be compiled and linked against the Orbiter library. 
The command 

\sourcecode{F_17_NTT_compile}

can be used to compile the code and run it. 
Note the dependency on the file \verb'NTT_k4_q17.cpp'. 
This means that \verb'make' would automatically 
invoke the first command if only the second one was issued.



\bigskip







